\documentclass[10pt]{article}

% ---------- Document Identity (update these only) ----------
\newcommand{\DocStage}{}
\newcommand{\DocTitle}{Your Road to Tier One SOF}
\newcommand{\ProgramName}{182-Week Civilian-to-Selection Readiness Pipeline}
\newcommand{\ProgramSubtitle}{}

% ---------- Page ----------
\usepackage[legalpaper,left=0.5in,right=0.5in,top=0.6in,bottom=0.45in]{geometry}

% ---------- Typography ----------
\usepackage[T1]{fontenc}
\usepackage[utf8]{inputenc}
\usepackage{libertinus}
\usepackage{microtype}
\usepackage{parskip}
\usepackage{ragged2e}
\usepackage{textcomp}
\AtBeginDocument{\color{black}}
\AtBeginDocument{\pagecolor{white}}
\AtBeginDocument{\justifying}

% ---------- Landscape pages ----------
\usepackage{pdflscape}

% ---------- Color ----------
\usepackage[table]{xcolor}
\definecolor{StageBlue}{HTML}{1F4E79}
\definecolor{StageBlueDark}{HTML}{163A5A}
\definecolor{LightGray}{HTML}{F2F4F7}
\definecolor{MidGray}{HTML}{D9DEE5}
\definecolor{RuleGray}{HTML}{C7CED8}
\definecolor{HeaderInk}{HTML}{000000}
\definecolor{Ink}{HTML}{000000}

% ---------- Utility ----------
\usepackage{enumitem}
\sloppy
\setlength{\emergencystretch}{2em}

% ---------- Boxes / UI ----------
\usepackage[most]{tcolorbox}

\tcbset{
  charterTitle/.style={
    enhanced,
    colback=StageBlue!4,
    colframe=StageBlue,
    boxrule=1.25pt,
    arc=3pt,
    left=16pt,right=16pt,top=14pt,bottom=14pt,
    borderline north={3.2pt}{0pt}{StageBlue},
    borderline west={4pt}{0pt}{StageBlue},
    borderline south={1.25pt}{0pt}{StageBlue},
    drop shadow={black!25!white},
  },
  charterRule/.style={
    enhanced,
    colback=LightGray,
    colframe=RuleGray,
    boxrule=0.85pt,
    arc=3pt,
    left=12pt,right=12pt,top=10pt,bottom=10pt,
    borderline west={3pt}{0pt}{StageBlue},
  },
  charterBadge/.style={
    enhanced,
    colback=StageBlue,
    coltext=white,
    colframe=StageBlueDark,
    boxrule=0.6pt,
    arc=2pt,
    left=6pt,right=6pt,top=3pt,bottom=3pt,
  }
}
\newtcbox{\Badge}{nobeforeafter,charterBadge}

% ---------- Lists ----------
\setlist[itemize]{leftmargin=1.5em, itemsep=0.25em, topsep=0.25em}
\setlist[enumerate]{leftmargin=1.8em, itemsep=0.25em, topsep=0.25em}

% ---------- TOC ----------
\usepackage{tocloft}
\setcounter{tocdepth}{3}
\setlength{\cftbeforesecskip}{3pt}
\setlength{\cftbeforesubsecskip}{2pt}
\setlength{\cftbeforesubsubsecskip}{0pt}
\renewcommand{\cftsecfont}{\bfseries}
\renewcommand{\cftsecpagefont}{\bfseries}
\renewcommand{\cftsubsecfont}{\normalfont}
\renewcommand{\cftsubsecpagefont}{\normalfont}
\renewcommand{\cftsubsubsecfont}{\normalfont}
\renewcommand{\cftsubsubsecpagefont}{\normalfont}
\setlength{\cftsecindent}{0pt}
\setlength{\cftsubsecindent}{1.25em}
\setlength{\cftsubsubsecindent}{2.8em}
\setlength{\cftsecnumwidth}{2.1em}
\setlength{\cftsubsecnumwidth}{2.8em}
\setlength{\cftsubsubsecnumwidth}{3.6em}

% Remove the built-in TOC heading so only the custom heading is shown
\makeatletter
\renewcommand{\tableofcontents}{\@starttoc{toc}}
\makeatother

% ---------- Header/Footer: page number only ----------
\usepackage{fancyhdr}
\pagestyle{fancy}
\fancyhf{}
\renewcommand{\headrulewidth}{0pt}
\renewcommand{\footrulewidth}{0pt}
\fancyfoot[C]{\thepage}

% ---------- Section formatting ----------
\usepackage{titlesec}

\titleformat{\section}
  {\Large\bfseries\color{Ink}}
  {\thesection}{0.6em}{}
  [\vspace{0.25em}\textcolor{StageBlue}{\rule{\linewidth}{0.8pt}}]

\titleformat{\subsection}
  {\large\bfseries\color{Ink}}
  {\thesubsection}{0.6em}{}

\titleformat{\subsubsection}
  {\normalsize\bfseries\color{Ink}}
  {\thesubsubsection}{0.6em}{}

\titlespacing*{\section}{0pt}{0.95em}{0.35em}
\titlespacing*{\subsection}{0pt}{0.65em}{0.25em}
\titlespacing*{\subsubsection}{0pt}{0.55em}{0.2em}

% ---------- Table engine ----------
\usepackage{tabularray}
\UseTblrLibrary{booktabs}

% ---------- tabularray: GLOBAL table treatment ----------
\SetTblrInner{
  cells = {fg=black, bg=white, valign=t},
}

% ---------- Header helpers ----------
\newcommand{\Hdr}[1]{\shortstack[c]{\strut #1\strut}}

% ---------- Lines ----------
\newcommand{\TitleLine}{\textcolor{StageBlue}{\rule{0.98\linewidth}{0.95pt}}}
\newcommand{\SubTitleLine}{\textcolor{RuleGray}{\rule{0.98\linewidth}{0.65pt}}}

% ---------- Continued on next page ----------
\newcommand{\ContinuedOnNextPageInline}{%
  \par
  \vfill
  \noindent\textcolor{RuleGray}{\rule{\linewidth}{1pt}}\vspace{-2pt}
  \noindent\hfill\textit{Continued on next page}\par\vspace{-10pt}
  \noindent\textcolor{RuleGray}{\rule{\linewidth}{1pt}}
  \clearpage
}

% ---------- PDF hyperlinks ----------
\usepackage[hidelinks]{hyperref}
\hypersetup{
  colorlinks=false,
  pdfborder={0 0 0},
  linktoc=all,
  pdftitle={\DocTitle},
  pdfauthor={\ProgramName}
}

% =========================================================
\begin{document}

\section*{Stage 2.E — Data Recording and Test Logistics Conventions}
\phantomsection
\addcontentsline{toc}{section}{Stage 2.E — Data Recording and Test Logistics Conventions}

\subsection*{2.E.1 Core Principles for Auditability and Comparability}
\phantomsection
\addcontentsline{toc}{subsection}{2.E.1 Core Principles for Auditability and Comparability}

\begin{itemize}
\item The PDF binder \textbf{MUST} be the single source of truth for decision-grade evidence used for gates. If evidence is not binder-recorded or binder-referenced, it is not admissible for gate decisions.
\item Runtime notes MAY exist outside the binder, but any gate claim \textbf{MUST} be supported by binder-recorded rows or binder-referenced Evidence Ref pointers that are stable and retrievable.
\item Every log row \textbf{MUST} be written with a no-inference posture: unknown values \textbf{MUST} be recorded as \textbf{UNKNOWN} or \textbf{MISSING}; they \textbf{MUST} NOT be guessed, interpolated, or reconstructed from memory for gate purposes.
\item Completeness of required fields is an admissibility control. If a required field is missing, the record is inadmissible for gates and \textbf{MUST} be treated as \textbf{INVALID} for gate use.
\item Identifiers (Metric \textbf{ID}, Protocol Card \textbf{ID}, Domain Tag, Environment Unit \textbf{ID}) \textbf{MUST} be treated as crosswalk keys. Key fields \textbf{MUST} remain stable; renamed or drifting keys break comparability and render downstream artifacts non-deployable.
\item Environment and tool stability is a comparability control. The environment and tool used for a test \textbf{MUST} be identifiable and stable enough to compare across time; material changes require new Environment Unit IDs or explicit tool-change records.
\item Deviations \textbf{MUST} be documented explicitly. If execution conditions differ from standard, the deviation \textbf{MUST} be recorded in the Notes field and, where applicable, the correct reason-code field (\textbf{SESSION-RC} or \textbf{TEST-RC}) \textbf{MUST} be populated.
\item Every test event log row \textbf{MUST} include Evidence Ref when required by the governing protocol card or logging rule. When artifact evidence is optional, Evidence Ref remains recommended but is not mandatory.
\item Replication status \textbf{MUST} be tracked for each Gate-Critical metric within the decision window using the required 1/3, 2/3, 3/3 notation. Tier classification requires replication per governance rules.
\item Timestamps \textbf{MUST} be decision-grade. All rows \textbf{MUST} use ISO 8601 timestamp format; timezone offset \textbf{MUST} be included when available. Missing timestamp or unverified late reconstruction makes the record inadmissible for gate decisions.
\item Reason-code namespaces \textbf{MUST} remain separated. \textbf{SESSION-RC} is used only for session \textbf{MODIFIED}/\textbf{INVALID} tagging; \textbf{TEST-RC} is used only for test \textbf{INVALID} tagging. The non-applicable reason-code field \textbf{MUST} be blank or N/A.
\item Rounding and unit rules \textbf{MUST} be enforced as anti-drift controls. Rounding and unit discipline is part of comparability; unit drift invalidates comparisons.
\end{itemize}

Binding binder architecture (required sections; exact names):

\begin{enumerate}
\item Daily Logs
\item Session Logs
\item Test Event Logs
\item Environment Unit Library
\item Weekly Audit and Gate Review
\item Patches and Change Log
\end{enumerate}

\subsubsection*{Assumption Ledger}
\phantomsection

\begin{itemize}
\item None.
\end{itemize}

\subsubsection*{Reconciliation Notes}
\phantomsection

\begin{itemize}
\item Conflict: Upstream logging conventions elsewhere include an intensity tier model as a required field. Adopted posture for Stage 2.E is explicitly locked by this prompt: Cardio intensity is recorded as RPE 0–10 plus Talk Test cue as the required controlled fields in Stage 2.E; any tier notation recorded elsewhere is supplemental only. Impact: Session Log schema in 2.E.4 records Cardio RPE, Talk Test cue, and Session RPE in distinct controlled fields; any tier notation, if recorded elsewhere, is treated as non-required free-text and is not a controlled field in Stage 2.E.
\end{itemize}

\newpage

\begin{landscape}

\subsection*{2.E.2 Required Identifiers}
\phantomsection
\addcontentsline{toc}{subsection}{2.E.2 Required Identifiers}

\begingroup
\scriptsize
\setlength{\tabcolsep}{2pt}
\renewcommand{\arraystretch}{1.12}

\begin{longtblr}{
  width=\linewidth,
  rowhead=1,
  colspec = {
    Q[l,wd=0.95in]
    X[l,1.15]
    X[l,1.35]
    X[l,2.55]
  },
  hlines,
  vlines,
  cells = {hsep=6pt, vsep=6pt, halign=l, valign=t},
  cell{1-Z}{1-4} = {fg=black},
  row{odd} = {bg=white},
  row{even} = {bg=LightGray},
  row{1} = {bg=StageBlue!15, fg=black, font=\bfseries, halign=l, valign=m},
}
\Hdr{Identifier} & \Hdr{Format} & \Hdr{Where Used} & \Hdr{Notes} \\
ISO Timestamp & YYYY-MM-DDThh:mm:ss±hh:mm & Daily Logs; Session Logs; Test Event Logs; Weekly Audit entries & If timezone offset is available, it \textbf{MUST} be recorded. If offset cannot be captured, the row \textbf{MUST} record ISO timestamp in UTC with “Z” only when a verifiable UTC source is used; otherwise ISO Timestamp is \textbf{MISSING} and the row is inadmissible for gates. \\
Phase \textbf{ID} & Phase 00–Phase 13 & Daily Logs; Session Logs; Test Event Logs; Weekly Audit and Gate Review & Phase numbering is canonical and \textbf{MUST} NOT be renumbered. \\
Phase Week & Wk-## (phase-relative) & Daily Logs; Session Logs; Test Event Logs; Weekly Audit and Gate Review & Phase Week is a phase-relative week index. If unknown, record \textbf{MISSING}. \\
Session Number & 1–6 & Session Logs & Exactly 6 sessions per week; no two-a-days. \\
Metric \textbf{ID} & \textbf{MET-2B}-### & Test Event Logs; Weekly Audit and Gate Review; Crosswalk hooks & Must match Stage 2.B. Undefined IDs are prohibited for deployable work. \\
Protocol Card \textbf{ID} & \textbf{C1}–\textbf{C18} & Test Event Logs & Must match Stage 2.C. If a test has no protocol, it is not gate-admissible. \\
Domain Tag & Stage 2.A canonical tags only & Test Event Logs; Weekly Audit and Gate Review & Must be one canonical tag per metric/test; no invented tags. \\
Session Validity Tag & \textbf{VALID} / \textbf{MODIFIED} / \textbf{INVALID} & Session Logs; Weekly Audit and Gate Review & Must match Stage 1.F. Missing required fields force \textbf{INVALID} for gate use. \\
Test Validity Tag & \textbf{VALID} / \textbf{INVALID} & Test Event Logs; Weekly Audit and Gate Review & Must match Stage 1.F. Missing required fields force \textbf{INVALID} for gate use. \\
\textbf{SESSION-RC} & \textbf{SESSION-RC}-### & Session Logs & Populate only when Session Validity Tag is \textbf{MODIFIED} or \textbf{INVALID}. If tag is \textbf{VALID}, \textbf{SESSION-RC} must be blank/N/A. \\
\textbf{TEST-RC} & \textbf{TEST-RC}-### & Test Event Logs & Populate only when Test Validity Tag is \textbf{INVALID}. If tag is \textbf{VALID}, \textbf{TEST-RC} must be blank/N/A. \\
Environment Unit \textbf{ID} & \textbf{ROUTE}-2E-### / \textbf{POOL}-2E-### / \textbf{TM}-2E-### / ROW-2E-### / BIKE-2E-### / ELL-2E-### & Test Event Logs; Environment Unit Library & Environment Unit \textbf{ID} is mandatory for run/ruck/swim/machine-based tests and for any test where environment comparability matters. \\
Evidence Ref & file name / URL / binder page ref / storage path pointer & Test Event Logs; Weekly Audit and Gate Review & Mandatory for every test event. Missing Evidence Ref makes the test inadmissible for gates. \\
Replication Status & 1/3, 2/3, 3/3 & Test Event Logs; Weekly Audit and Gate Review & Must reflect count of \textbf{VALID} events for that metric within the current decision window. \\
Tool Standard (Descriptor) & Free-text descriptor of timer, scale, or tape & Test Event Logs & Tool IDs are optional; when no Tool \textbf{ID} exists, descriptor is required (e.g., “Phone Timer,” “Digital Scale,” “Cloth Tape”). \\
\end{longtblr}

\endgroup

\end{landscape}

\begin{landscape}

\subsection*{2.E.3 Daily Log Minimum Dataset}
\phantomsection
\addcontentsline{toc}{subsection}{2.E.3 Daily Log Minimum Dataset}

\begingroup
\scriptsize
\setlength{\tabcolsep}{2pt}
\renewcommand{\arraystretch}{1.12}

\begin{longtblr}{
  width=\linewidth,
  rowhead=1,
  colspec = {
    Q[l,wd=0.90in]
    Q[l,wd=0.90in]
    Q[l,wd=0.70in]
    Q[l,wd=0.90in]
    Q[l,wd=0.90in]
    Q[l,wd=0.90in]
    Q[l,wd=1.25in]
    Q[l,wd=0.80in]
    X[l,1.30]
  },
  hlines,
  vlines,
  cells = {hsep=6pt, vsep=6pt, halign=l, valign=t},
  cell{1-Z}{1-9} = {fg=black},
  row{odd} = {bg=white},
  row{even} = {bg=LightGray},
  row{1} = {bg=StageBlue!15, fg=black, font=\bfseries, halign=l, valign=m},
}
\Hdr{ISO Timestamp} &
\Hdr{Phase \textbf{ID} / Phase Week} &
\Hdr{Sleep Hrs} &
\Hdr{Sleep Qual 0–10} &
\Hdr{Soreness 0–10} &
\Hdr{Fatigue 0–10} &
\Hdr{Pain 0–10 plus location} &
\Hdr{Steps count} &
\Hdr{Notes} \\
\end{longtblr}

\endgroup

Rules (binding):

\begin{itemize}
\item Steps count is required daily.
\item Pain location is required when pain is greater than 0.
\item Missing fields \textbf{MUST} be recorded as \textbf{MISSING} (or \textbf{UNKNOWN} where appropriate). Values \textbf{MUST} NOT be inferred or backfilled from memory for gate decisions.
\item If ISO Timestamp is \textbf{MISSING}, the daily log row is non-auditable and \textbf{MUST} NOT be used in gate reasoning.
\end{itemize}

\subsection*{2.E.4 Session Log Minimum Dataset}
\phantomsection
\addcontentsline{toc}{subsection}{2.E.4 Session Log Minimum Dataset}

\begingroup
\scriptsize
\setlength{\tabcolsep}{2pt}
\renewcommand{\arraystretch}{1.12}

\begin{longtblr}{
  width=\linewidth,
  rowhead=1,
  colspec = {
    Q[l,wd=.65in]
    Q[l,wd=0.575in]
    Q[l,wd=0.3in]
    Q[l,wd=.50in]
    Q[l,wd=0.50in]
    Q[l,wd=.45in]
    Q[l,wd=.65in]
    Q[l,wd=0.70in]
    Q[l,wd=0.35in]
    Q[l,wd=0.45in]
    Q[l,wd=.55in]
    Q[l,wd=0.45in]
    Q[l,wd=0.60in]
    Q[l,wd=0.4in]
    Q[l,wd=0.4in]
    X[l,1.10]
  },
  hlines,
  vlines,
  cells = {hsep=6pt, vsep=6pt, halign=l, valign=t},
  cell{1-Z}{1-16} = {fg=black},
  row{odd} = {bg=white},
  row{even} = {bg=LightGray},
  row{1} = {bg=StageBlue!15, fg=black, font=\bfseries, halign=l, valign=m},
}
\Hdr{ISO Timestamp} &
\Hdr{Phase \textbf{ID} / Phase Week} &
\Hdr{Session Number} &
\Hdr{Validity Tag} &
\Hdr{\textbf{SESSION-RC}} &
\Hdr{Cardio Duration} &
\Hdr{Tool Standard} &
\Hdr{Cardio Modality} &
\Hdr{Cardio RPE} &
\Hdr{Talk Test Cue} &
\Hdr{Main Work} &
\Hdr{Session RPE} &
\Hdr{Cardio Minimum Standard Met Y/N} &
\Hdr{Pain Flags} &
\Hdr{Foot Flags} &
\Hdr{Notes} \\
\end{longtblr}

\endgroup

Rules (binding):

\begin{itemize}
\item Session Validity Tag \textbf{MUST} be one of \textbf{VALID} / \textbf{MODIFIED} / \textbf{INVALID}.
\item \textbf{SESSION-RC} is populated only when Session Validity Tag is \textbf{MODIFIED} or \textbf{INVALID}. If Session Validity Tag is \textbf{VALID}, \textbf{SESSION-RC} \textbf{MUST} be blank or N/A.
\item Cardio Min Standard Met Y/N is a required field and \textbf{MUST} be derived from the session record (Cardio duration, continuity, and documentation completeness).
\item Cardio requirements \textbf{MUST} be expressed in body text only and \textbf{MUST} NOT be embedded into headings or column headers.
\item Cardio Intensity \textbf{MUST} be recorded as RPE 0–10 plus a Talk Test cue (no tier model in Stage 2.E).
\item If any required field is \textbf{MISSING}, the session evidence is inadmissible for gate decisions and \textbf{MUST} be treated as \textbf{INVALID} for gate use.
\end{itemize}

\subsection*{2.E.5 Test Event Log Minimum Dataset}
\phantomsection
\addcontentsline{toc}{subsection}{2.E.5 Test Event Log Minimum Dataset}

\begingroup
\scriptsize
\setlength{\tabcolsep}{2pt}
\renewcommand{\arraystretch}{1.12}

\begin{longtblr}{
  width=\linewidth,
  rowhead=1,
  colspec = {
    Q[l,wd=.435in]
    Q[l,wd=0.285in]
    Q[l,wd=0.40in]
    Q[l,wd=.50in]
    Q[l,wd=0.50in]
    Q[l,wd=.60in]
    Q[l,wd=.75in]
    Q[l,wd=0.85in]
    Q[l,wd=0.50in]
    Q[l,wd=0.35in]
    Q[l,wd=0.30in]
    Q[l,wd=.35in]
    Q[l,wd=0.45in]
    X[l,1.25]
    X[l,0.80]
  },
  hlines,
  vlines,
  cells = {hsep=6pt, vsep=6pt, halign=l, valign=t},
  cell{1-Z}{1-15} = {fg=black},
  row{odd} = {bg=white},
  row{even} = {bg=LightGray},
  row{1} = {bg=StageBlue!15, fg=black, font=\bfseries, halign=l, valign=m},
}
\Hdr{ISO\\ Timestamp} &
\Hdr{Phase \\\& Week} &
\Hdr{Protocol \\Card \textbf{ID}} &
\Hdr{Linked \\Metric \textbf{ID}(s)} &
\Hdr{Domain\\ Tag} &
\Hdr{Environment \\Unit \textbf{ID}} &
\Hdr{Tool Standard (Timer / Scale / Tape)} &
\Hdr{Warm-Up \\Completed Y/N} &
\Hdr{Test \\Validity Tag} &
\Hdr{\textbf{TEST}-\\RC} &
\Hdr{Result} &
\Hdr{Evidence\\ Ref} &
\Hdr{Replication \\Status} &
\Hdr{Conditions\\ Notes} &
\Hdr{Stop/Abort} \\
\end{longtblr}

\endgroup

\newpage

Rules (binding):

\begin{itemize}
\item Test Validity Tag \textbf{MUST} be \textbf{VALID} or \textbf{INVALID}.
\item \textbf{TEST-RC} is populated only when Test Validity Tag is \textbf{INVALID}. If Test Validity Tag is \textbf{VALID}, \textbf{TEST-RC} \textbf{MUST} be blank or N/A.
\item Evidence Ref is mandatory. Missing Evidence Ref makes the test inadmissible for gate decisions and \textbf{MUST} be treated as \textbf{INVALID} for gate use.
\item Environment Unit \textbf{ID} is mandatory for any test where environment comparability is relevant (\textbf{RUN}, \textbf{RUCK}, \textbf{SWIM}, \textbf{AER} machine-based tests). Missing Environment Unit \textbf{ID} makes the test inadmissible for gate decisions and \textbf{MUST} be treated as \textbf{INVALID} for gate use.
\item If one protocol produces multiple metrics, Linked Metric \textbf{ID}(s) \textbf{MUST} list all \textbf{MET-2B}-### IDs separated by semicolons.
\item Stop Abort is a required field. If Yes, record the observable reason and the immediate action taken (termination posture and escalation when indicated).
\end{itemize}

\subsection*{2.E.6 Environment Recording Standards for Route Pool and Machine}
\phantomsection
\addcontentsline{toc}{subsection}{2.E.6 Environment Recording Standards for Route Pool and Machine}

Required environment recording posture (binding):

\begin{itemize}
\item Every \textbf{RUN}, \textbf{RUCK}, \textbf{SWIM}, and machine-based \textbf{AER} test \textbf{MUST} reference a registered Environment Unit \textbf{ID} from the Environment Unit Library.
\item If the environment changes materially (route distance verification method changes; pool unit/length changes; treadmill make/model changes; lane conditions materially change; surface changes materially), a new Environment Unit \textbf{ID} \textbf{MUST} be created. Prior IDs \textbf{MUST} NOT be reused to represent a different environment state.
\item Environment Unit IDs are crosswalk keys. They \textbf{MUST} be stable, unique, and non-reused.
\end{itemize}

Naming structures (binding formatting conventions):

\begin{itemize}
\item Route \textbf{ID} structure: \textbf{ROUTE}-2E-### — distance label — surface — location note
\item Pool \textbf{ID} structure: \textbf{POOL}-2E-### — length and units — location note
\item Machine \textbf{ID} structures: \textbf{TM}-2E-###, ROW-2E-###, BIKE-2E-###, ELL-2E-###
\end{itemize}

Environment Unit Library registry table schema (exact columns):

\begingroup
\scriptsize
\setlength{\tabcolsep}{2pt}
\renewcommand{\arraystretch}{1.12}

\begin{longtblr}{
  width=\linewidth,
  rowhead=1,
  colspec = {
    Q[l,wd=1.00in]
    Q[l,wd=1.10in]
    Q[l,wd=1.25in]
    Q[l,wd=.50in]
    X[l,1.00]
    Q[l,wd=1.60in]
    Q[l,wd=.875in]
    X[l,0.80]
  },
  hlines,
  vlines,
  cells = {hsep=6pt, vsep=6pt, halign=l, valign=t},
  cell{1-Z}{1-8} = {fg=black},
  row{odd} = {bg=white},
  row{even} = {bg=LightGray},
  row{1} = {bg=StageBlue!15, fg=black, font=\bfseries, halign=l, valign=m},
}
\Hdr{Environment Unit \textbf{ID}} &
\Hdr{Type (Route / Pool / Machine)} &
\Hdr{Standard Distance or Length} &
\Hdr{Location} &
\Hdr{Surface or Conditions} &
\Hdr{Grade or Incline} &
\Hdr{Verification Method} &
\Hdr{Notes} \\
\end{longtblr}

\endgroup

Required condition fields to capture (record in Notes or Surface/Conditions as applicable):

\begin{itemize}
\item Outdoor routes: weather category; surface condition; wind note if material; temperature if known; traffic/hazard notes when relevant.
\item Pools: length units (yards or meters); lane condition; crowding note; water temperature if known; facility rule constraints when relevant.
\item Machines: make and model; console distance basis (displayed distance); grade/incline and setting capture expectations per test.
\end{itemize}

\end{landscape}

\subsection*{2.E.7 Tool Standardization Rules}
\phantomsection
\addcontentsline{toc}{subsection}{2.E.7 Tool Standardization Rules}

Timer standardization posture (binding):

\begin{itemize}
\item The timing method used for time trials \textbf{MUST} be recorded for every test event (Tool Standard field).
\item A timer change event \textbf{MUST} be recorded when the timing device or method changes (e.g., phone timer to watch timer; different app; different watch).
\item If timing method is not recorded, the test is \textbf{INVALID} for gate use.
\end{itemize}

Scale standardization posture (binding):

\begin{itemize}
\item Bodyweight \textbf{MUST} be rounded to the nearest 0.1 lb.
\item Time-of-day and measurement conditions \textbf{SHOULD} be consistent within trend windows; if they change, the change \textbf{MUST} be recorded in \textbf{MEASURE} or Notes.
\item Scale identity \textbf{MUST} be recorded (Tool \textbf{ID} if it exists, otherwise a clear description or stable label).
\item Scale change events \textbf{MUST} be recorded. A scale change is a comparability break until stabilized and documented.
\end{itemize}

Tape standardization posture (binding):

\begin{itemize}
\item Tape sites \textbf{MUST} be rounded to the nearest 0.25 in.
\item Landmark definitions and side selection \textbf{MUST} be stable within a trend window; any change is a comparability break and \textbf{MUST} be recorded.
\item Tape identity \textbf{MUST} be recorded (Tool \textbf{ID} if it exists, otherwise a clear description or stable label).
\end{itemize}

Rounding and unit rules (binding anti-drift):

\begin{itemize}
\item Bodyweight: nearest 0.1 lb.
\item Tape sites: nearest 0.25 in.
\item Time trials: nearest 1 second.
\item Distances: record the named test distance (e.g., “1.5 mi,” “500 yd,” “500 m,” “3 mi”) and \textbf{MUST} NOT overwrite with GPS-rounded values in the Result field.
\end{itemize}

Tool change event definition (binding):

\begin{itemize}
\item A tool change event occurs when any primary measurement tool changes in a way that could affect comparability: timer device or method, scale device, tape measure, pool length or unit, route-verification method, treadmill or machine make/model.
\item Tool change events \textbf{MUST} be logged in Notes and, when relevant, in the Environment Unit Library (new Environment Unit \textbf{ID} when the environment/tool is part of the environment definition).
\end{itemize}

Tool drift and conservative handling (binding):

\begin{itemize}
\item When tool drift affects comparability (different treadmill, different pool unit/length, different route verification method), results MAY be recorded, but \textbf{MUST} be interpreted conservatively for gate use until stable, replicated \textbf{VALID} evidence exists under the new stable environment/tool state.
\item If tool drift prevents validity condition requirements from being met (missing identity, missing verification method), the result is \textbf{INVALID} for gate use.
\end{itemize}

\newpage

\subsection*{2.E.8 Missing Data and Deviation Handling}
\phantomsection
\addcontentsline{toc}{subsection}{2.E.8 Missing Data and Deviation Handling}

No inference rule (binding):

\begin{itemize}
\item Unknown values \textbf{MUST} be recorded as \textbf{UNKNOWN} or \textbf{MISSING}. No inference is permitted. Do not estimate, interpolate, or “fill in” for completeness.
\end{itemize}

Deviation documentation rule (binding):

\begin{itemize}
\item Any deviation from standard conditions \textbf{MUST} be documented with:
\begin{itemize}
\item the applicable reason code field (\textbf{SESSION-RC} for sessions; \textbf{TEST-RC} for tests) when the validity tag is \textbf{MODIFIED}/\textbf{INVALID} (session) or \textbf{INVALID} (test), and
\item a free-text note describing what changed, why, and the observed impact on validity/admissibility posture.
\end{itemize}
\end{itemize}

Missing required fields rule (binding):

\begin{itemize}
\item If a session or test was performed but required fields are missing, the evidence is inadmissible for gate decisions and \textbf{MUST} be treated as \textbf{INVALID} for gate use automatically.
\item Missing required fields \textbf{MUST} NOT be repaired by memory for gate purposes.
\end{itemize}

Late entry rule (binding):

\begin{itemize}
\item Late entry is permitted only when the missing field can be verified with objective evidence (e.g., a timestamped photo/video, device log export, or facility record).
\item A late entry \textbf{MUST} be marked explicitly in Notes as \textbf{VERIFIED} \textbf{LATE} \textbf{ENTRY} and \textbf{MUST} include the objective evidence pointer in Evidence Ref (tests) or Notes (sessions).
\item If objective verification is not available, the entry remains inadmissible for gate decisions.
\end{itemize}

Retest posture after invalid test (binding):

\begin{itemize}
\item \textbf{INVALID} tests provide no gate credit and \textbf{MUST} be repeated under controlled conditions if the test is gate-required.
\item Retesting \textbf{MUST} prioritize standardization and safety; retesting \textbf{MUST} NOT become PR chasing or justification for unsafe escalation.
\end{itemize}

Environment change rule (binding):

\begin{itemize}
\item When an environment changes materially, create a new Environment Unit \textbf{ID} rather than altering prior records.
\item Do not rewrite history. Prior test events remain interpreted under the environment/tool state recorded at the time.
\end{itemize}

\newpage

\subsection*{2.E.9 Weekly Compliance and Audit Checks}
\phantomsection
\addcontentsline{toc}{subsection}{2.E.9 Weekly Compliance and Audit Checks}

Weekly Audit and Gate Review checklist (binding; numeric defaults explicit):

\begin{itemize}
\item Verify admissible session compliance: at least 5 of 6 admissible sessions per week on average across the decision window.
\item Verify \textbf{MODIFIED} cap posture: no more than 1 \textbf{MODIFIED} session per week outside medical issues; document medical issue justification when exceeded.
\item Verify Cardio minimum standard met rate across the decision window is checked and recorded (and any failures are visible).
\item Review session validity distribution (counts of \textbf{VALID} / \textbf{MODIFIED} / \textbf{INVALID}) and document dominant failure causes.
\item Review test validity rate (\textbf{VALID} vs \textbf{INVALID}) for any tests executed in or affecting the decision window.
\item Review Replication Status (1/3, 2/3, 3/3) for Gate-Critical metrics tested; do not claim tier classification without replication.
\item Confirm Environment Unit IDs are stable and any changes are logged with new IDs where required.
\item Confirm tool changes are logged (timer, scale, tape, or device changes) and comparability breaks are documented.
\item Review pain and foot flags for trend escalation; document constraints and conservative actions.
\item Enforce advancement posture constraint: gait-altering pain flag must be No for advancement posture.
\item Review Stop/Abort events; verify documentation includes observable trigger, immediate action, escalation path when indicated, and disposition.
\item Confirm Evidence Ref completeness for all test event rows where Evidence Ref is required by protocol or logging doctrine; missing required Evidence Ref is an automatic gate inadmissibility trigger for that test.
\end{itemize}

\subsection*{2.E.10 Conflict-Resolution Rule}
\phantomsection
\addcontentsline{toc}{subsection}{2.E.10 Conflict-Resolution Rule}

\begin{itemize}
\item Any session or test lacking required logging fields is inadmissible for gate decisions and \textbf{MUST} be treated as \textbf{INVALID} for gate use.
\item If any downstream artifact attempts to use inadmissible evidence to justify thresholds, tier classification, or gate outcomes, that artifact is Draft — Non-Deployable until corrected and re-versioned.
\end{itemize}


\end{document}